% ECONOMICS_PROBLEM: {"id": "H22-468", "title": "Tax incidence anomalies", "jel_code": "H22", "source_id": "OP-0476"}
% !TeX program = xelatex
\documentclass[11pt,a4paper]{article}
\usepackage[margin=23mm]{geometry}
\usepackage{fontspec}
\setmainfont{Latin Modern Roman}
\usepackage{amsmath,amssymb,mathtools}
\usepackage{xcolor,enumitem,booktabs,longtable,array,xurl,hyperref}
\definecolor{muted}{HTML}{53636C}
\hypersetup{colorlinks=true,urlcolor=blue,linkcolor=blue}
\setlength{\parindent}{0pt}
\setlength{\parskip}{5pt}
\newcommand{\R}{\mathbb R}
\newcommand{\E}{\mathbb E}
\newcommand{\N}{\mathbb N}
\newcommand{\ind}{\mathbf 1}
\newcommand{\Var}{\operatorname{Var}}
\newcommand{\Cov}{\operatorname{Cov}}
\newcommand{\supp}{\operatorname{supp}}
\DeclareMathOperator*{\argmin}{arg\,min}
\DeclareMathOperator*{\argmax}{arg\,max}
\newcommand{\entrymeta}[3]{{\sffamily\small\color{muted}\textbf{\texttt{#1}}\quad|\quad #2\par #3\par}}
\newcommand{\Problemref}[1]{\texttt{#1}}
\newcommand{\JELref}[2]{\texttt{#2} (JEL \texttt{#1})}
\begin{document}
\section*{Tax incidence anomalies}
\textbf{Problem ID:} \texttt{H22-468}\quad \textbf{JEL:} \texttt{H22}\par
% BEGIN ECONOMICS STATEMENT
\label{problem:OP-0476}
{\sffamily\small\color{muted}Original subject group: Taxation and social insurance\par}
\label{definitions:problem:30007758}
\textbf{Audit classification:} Background; proposed extension. The Annual Reviews record verifies title, author, year, pages, DOI and anomaly categories. Its accessible abstract is background for the exact directional-derivative problem. Corrected conflation of finite asymmetry, differentiability, and hysteresis.
\subsection*{Definitions and precise research target}
Fix a market, common preannouncement state \(z\), price aggregation weights, policy duration and a tax unit. Use a specific tax \(\delta\) in currency per unit; for an ad valorem rate use separate normalization. Let \(p_h(\delta;z)\) be the mean tax-inclusive unit price at horizon \(h\) under an intervention changing the same baseline tax by \(\delta\), with fixed information at announcement and declared equilibrium selection. If finite one-sided derivatives exist, set
\[
\begin{aligned}
 \rho_+(h,z)&=\lim_{\delta\downarrow0}\frac{p_h(\delta;z)-p_h(0;z)}{\delta},\\
 \rho_-(h,z)&=\lim_{\delta\uparrow0}\frac{p_h(\delta;z)-p_h(0;z)}{\delta},\\
 A(h,z)&=\rho_+(h,z)-\rho_-(h,z).
\end{aligned}
\]
If \(p_h\) is differentiable at zero, \(A(h,z)=0\) by definition. Unequal responses to two finite reforms therefore do not prove a derivative kink. For fixed \(d>0\), define a distinct finite-change target
\[
 A_d(h,z)=\frac{p_h(d;z)-p_h(0;z)}d
 -\frac{p_h(-d;z)-p_h(0;z)}{-d}.
\]
Nonzero \(A_d\) is possible under smooth nonlinear pricing. Permanent increases and temporary reductions have different intervention paths and cannot be substituted for these symmetric treatments. With observational law \(P_O\), fix an identification class covering reform assignment, anticipation, latent state and equilibrium restrictions; characterize the joint identified set of \(A\) and \(A_d\). A mechanism claim must specify how its primitives produce asymmetry at a common initial state and give extra observable restrictions or interventions that distinguish it from alternatives. For example, nonsmooth price adjustment can produce a kink, while history-dependent states can produce persistent level differences without a common-state kink. This precise directional/finite-change distinction is editorial; the cited review documents tax-incidence anomalies.
\subsection*{Variants and assumption changes}
\begin{enumerate}
\item \textbf{Infinitesimal versus finite reforms (editorial).} Estimate \(A\) near zero or \(A_d\) at a specified size. The first detects nondifferentiability; the second also detects ordinary curvature.
\item \textbf{Common state versus hysteresis (editorial).} Compare two fresh reforms from \(z\), or a tax rise followed by reversal along a changing state path. Persistent prices after reversal are a path-dependence target.
\item \textbf{Statutory remitter (editorial).} Hold the tax wedge fixed while changing whether buyers or sellers remit. A changed response then tests statutory-incidence invariance, a separate anomaly.
\end{enumerate}
\subsection*{References and source audit}
Benzarti (2025), Annual Review of Economics 17, 615--634, publisher abstract: metadata and anomaly categories verified. Full-text loading failed; the narrower exact mathematical question is not attributed to that abstract. The earlier NBER version is distinct.
\begin{enumerate}\raggedright
\item\hypertarget{definitions:source:30007758:1}{} Youssef Benzarti. 2025. \emph{Tax Incidence Anomalies}. Annual Review of Economics 17, 615–634; abstract.
\url{https://www.annualreviews.org/content/journals/10.1146/annurev-economics-081324-085805}.
DOI: \url{https://doi.org/10.1146/annurev-economics-081324-085805}.
\end{enumerate}

\subsection*{Common conventions: Source mathematical conventions}

These conventions apply to each entry unless explicitly replaced.

\textbf{Models and identification.} A primitive model \(m\) specifies all probability laws, feasible choices, information, equilibrium and measurement rules used by its target. The admissible class \(\mathcal M\) is fixed independently of outcomes. If \(P_O(m)\) is its observable probability law and \(\tau(m)\) the target, its identified set at observable law \(P\) is
\[
 \mathcal I_\tau(P)=\{\tau(m):m\in\mathcal M,\ P_O(m)=P\}.
\]
Point identification means that this set is a singleton. An empty set signals incompatibility of data and assumptions. Coordinate bounds need not describe the joint set. Bounds are sharp only if every retained target is attainable or arbitrarily approachable by compatible models. Finite bounds require explicit restrictions on unobserved quantities; unrestricted confounding, hidden wealth, extrapolation or arbitrary equilibrium selection can yield uninformative sets.

\textbf{Causal interventions.} A potential outcome \(Y_i(p)\) is the outcome under a complete policy regime \(p\), including timing, financing, information and market spillovers. The target population and units are fixed before treatment. With interference, use the complete assignment vector or a justified exposure mapping; individual no-interference notation alone cannot identify an economy-wide intervention. A difference of mean potential outcomes does not require the cross-world joint distribution. Conditioning on post-treatment migration, survival, enrollment or employment changes the estimand and needs separate justification.

\textbf{Statistical statements.} An estimator, confidence set, forecast or test specifies a sampling law, model class and loss/coverage criterion. Uniform \((1-\alpha)\)-coverage means \(\inf_{m\in\mathcal M}\Pr_m\{\tau(m)\in C_n\}\geq1-\alpha\), or an explicitly asymptotic version. The whole parameter space trivially has coverage, so informativeness requires a stated width, risk or power objective. A forecasting improvement is relative to a named benchmark, information set and evaluation population.

\textbf{Equilibrium and optimization.} An equilibrium comprises feasible plans, optimality, beliefs where needed, budget constraints and all market/resource clearing. The primitives or a stated benchmark define each correspondence \(\mathcal E(p;m)\). Nonemptiness is assumed or proved, never inferred just from compact policy sets. A selected-equilibrium target uses a declared selection rule; a robust target quantifies over all permitted equilibria. Use suprema or infima unless attainment is established. An \(\varepsilon\)-optimal policy is feasible and within \(\varepsilon\) of the finite optimum value. Report an empty feasible set separately.

\textbf{Welfare and accounting.} Normative utility, interpersonal normalization, social weights, numeraire, discounting and the use of public receipts are primitives. Behavioral utility need not coincide with the welfare criterion. Household welfare includes ownership income and rents once; adding profits again to compensating variation generally double-counts them. A monetary equivalent is defined by an explicit transfer and price convention, with monotonicity and range conditions. Average effects, mechanism contributions, transfers and welfare losses are different targets. Infinite sums require convergence and dynamic feasible plans require terminal or no-Ponzi conditions.

\textbf{Source versions.} A working-paper date, revision date and journal publication year are distinguished. A DOI for a journal version is not automatically the DOI of an earlier manuscript. Locators refer to the stated version and printed pages where specified. A metadata-only check does not verify a claim about a full-text conclusion. Every entry's source subsection narrows the provenance claim accordingly.
% END ECONOMICS STATEMENT
\end{document}
